\documentclass[12pt]{article}
\usepackage[T1]{fontenc}
\usepackage[utf8]{inputenc}
\usepackage{lmodern}
\usepackage{textcomp}
\usepackage{enumitem}
\usepackage{geometry}
\geometry{margin=1in}
\begin{document}
\begin{center}
Answer Key - Science - K-12 Level Assessment 19 \
\small (Answers are ordered exactly as in the question paper.)
\end{center}

\section*{Multiple Choice}
\begin{enumerate}[label=\arabic*.]
\item \textbf{Answer:} B
\\ \textit{Options:} A) 1 kg subjected to a force of 1 N; B) 1 kg subjected to a force of 100 N; C) 100 kg subjected to a force of 1 N; D) 100 kg subjected to a force of 100 N
\\ \textit{Note:} The correct answer is 1 kg subjected to a force of 100 N because, according to Newton's second law of motion (F=ma), the acceleration is directly proportional to the force applied and inversely proportional to the mass, resulting in the highest acceleration for this scenario.
\item \textbf{Answer:} A
\\ \textit{Options:} A) Rocks have many scratches.; B) Rocks contain fossils.; C) Rocks are covered by moss.; D) Rocks are deep underground.
\\ \textit{Note:} The presence of scratches on rocks, known as striations, is a key indicator of glacial movement, as they are formed when glaciers drag debris across the landscape, providing strong evidence that glaciers once covered the area.
\item \textbf{Answer:} D
\\ \textit{Options:} A) Earth.; B) the Moon.; C) the galaxy.; D) the Sun.
\\ \textit{Note:} The correct answer is "the Sun" because, due to the force of gravity, planets are attracted to the Sun, which acts as the central point in our solar system around which they orbit.
\item \textbf{Answer:} A
\\ \textit{Options:} A) rabbit, seed, bird; B) seed, bird, wind; C) Volcano, candle, rabbit; D) wind, candle, volcano
\\ \textit{Note:} The group "rabbit, seed, bird" consists only of living things because rabbits and birds are animals, while seeds are the reproductive structures of plants, all of which are classified as living organisms.
\item \textbf{Answer:} D
\\ \textit{Options:} A) a map showing the routes of city buses; B) a map showing the locations of streets; C) a map showing the locations of houses; D) a map showing the elevations of ground surfaces
\\ \textit{Note:} A map showing the elevations of ground surfaces is the best choice for analyzing potential flood areas because it allows the engineer to identify low-lying regions that are more likely to collect water during heavy rains.
\end{enumerate}

\section*{True / False}
\begin{enumerate}[label=\arabic*.]
\setcounter{enumi}{5}
\item \textbf{Answer:} True
\\ \textit{Options:} A) True; B) False
\\ \textit{Note:} A galaxy is composed of billions to trillions of stars, along with gas, dust, and dark matter, making it significantly larger in size compared to individual astronomical objects like black holes and neutron stars.
\item \textbf{Answer:} True
\\ \textit{Options:} A) True; B) False
\\ \textit{Note:} Preventing sperm production in male insects leads to reduced mating success, resulting in fewer fertilized eggs and ultimately lowering the overall insect population.
\item \textbf{Answer:} False
\\ \textit{Options:} A) True; B) False
\\ \textit{Note:} Pioneer plants are typically small, hardy species that are adapted to harsh conditions and can establish themselves quickly in disturbed environments without needing to grow large.
\item \textbf{Answer:} False
\\ \textit{Options:} A) True; B) False
\\ \textit{Note:} Reef ecosystems are typically found in stable, warm waters and do not experience the extreme temperature fluctuations that intertidal zones do, which are subject to significant changes due to tidal influences.
\item \textbf{Answer:} True
\\ \textit{Options:} A) True; B) False
\\ \textit{Note:} Enzymes lower the activation energy required for biochemical reactions, allowing them to proceed more efficiently at the lower temperatures typically found in living organisms.
\end{enumerate}

\section*{Long Form Response}
\begin{enumerate}[label=\arabic*.]
\setcounter{enumi}{10}
\item \textbf{Answer:} To determine which plant food, "Quickgrow" or "Supergrow," promotes the best growth in a specific type of houseplant, Pat can design a controlled experiment. First, he should select a uniform group of the same type of houseplant and divide them into two groups, ensuring they are all labeled clearly. Next, he should apply "Quickgrow" to one group and "Supergrow" to the other, while keeping all other conditions, such as light and water, consistent. Over a set period, Pat can measure and record the growth of the plants in height or number of leaves, comparing the results at the end of the experiment. This structured approach will help him identify which plant food is more effective in promoting growth.
\\ \textit{Note:} correct because it outlines a controlled experiment that involves the use of a uniform plant group, clear labeling, consistent environmental conditions, and systematic measurement of growth to effectively compare the effects of the two plant foods on growth.
\item \textbf{Answer:} Matter can be classified into different types based on its properties, including solids, liquids, and gases. Solids, like ice, have a definite shape and volume due to tightly packed particles, which means they maintain their form even when placed in different containers. Liquids, such as water, have a definite volume but take the shape of their container, allowing them to flow and adapt, which is essential for everyday tasks like drinking or washing. Gases, like oxygen, have neither a definite shape nor volume, expanding to fill any space, which is crucial for processes like breathing. Understanding these properties helps us apply matter effectively in various contexts, from building sturdy structures with solids to using gases in breathing and combustion.
\\ \textit{Note:} The answer correctly identifies and describes the three main states of matter-solids, liquids, and gases-by explaining their distinct properties and providing relevant everyday examples, thereby demonstrating how these characteristics influence their behavior and applications.
\end{enumerate}

\vfill
\noindent\textit{End of Answer Key}
\end{document}